% Guide 2 — objective: read a log-mel spectrogram with understanding
% (plan 019, "The second study-guide objective"; ADR 008).
% Build: make study   ->  air-to-log-mel.pdf (this) + the solutions manual.
\providecommand{\StudyRoot}{..}
\documentclass[11pt]{report}
\usepackage{theme}
\usepackage[hidelinks]{hyperref}
% Generated by tools/build_anchors.py from anchors.yaml — do not edit.
% Each macro splices a verified value; the cited plan anchor is the
% argument of record.

% 001.astronomical: order of C(T+U, T-U) at T=100, U=50 [deep_learning/README.md row 4 (plan 001 research pass)]
\expandafter\newcommand\csname anchor@001.astronomical\endcsname{2 \times 10^{40}}
% 001.paths15: paths collapsing to AB at T=4 (= C(6,2)) [docs/plans/001-ctc-alignment.md; verified two ways (enumeration + trellis)]
\expandafter\newcommand\csname anchor@001.paths15\endcsname{15}
% 001.raw81: raw path count, 3 symbols over 4 frames (3^4) [docs/plans/001-ctc-alignment.md; deep_learning/README.md row 4]
\expandafter\newcommand\csname anchor@001.raw81\endcsname{81}
% 009.G.nllgradient: gradient of the NLL at z = (2, 1, 0), correct class first [docs/plans/009-calculus-derivatives.md, anchor G]
\expandafter\newcommand\csname anchor@009.G.nllgradient\endcsname{(-0.3348,\ 0.2447,\ 0.0900)}
% 010.K.NLL: -ln P(AB|X), displayed at 4 dp [docs/plans/010-ctc-gradient.md, anchor K]
\expandafter\newcommand\csname anchor@010.K.NLL\endcsname{0.7181}
% 010.K.P: P(AB|X) on the worked matrix, exact (4877/10000) [docs/plans/010-ctc-gradient.md, anchor K]
\expandafter\newcommand\csname anchor@010.K.P\endcsname{0.4877}
% 010.L.uniformP: P(AB|X) under uniform outputs, 15 * (1/3)^4 [docs/plans/010-ctc-gradient.md, anchor L]
\expandafter\newcommand\csname anchor@010.L.uniformP\endcsname{5/27}
% 010.N.t5paths: paths for a 2-letter transcript at T=5 (C(7,3)) [docs/plans/010-ctc-gradient.md, anchor N]
\expandafter\newcommand\csname anchor@010.N.t5paths\endcsname{35}


% Guide-owned assembly macros (plan 012 R4) — split from guide.tex so
% tools/check_study_layout.py can render macro-generated figures
% standalone.
% The roadmap — guide-owned assembly state. \roadmap{n} draws the
% objective's four stations with the first n filled (read). Stations past
% \builtstations are drawn dashed: planned, not yet written — the picture
% must not claim chapters the guide does not contain. Not shared with
% ctc-algorithm/macros.tex: that roadmap has no planned state, and a
% guide's road is its own.
\newcommand{\roadmapstations}{%
  the spectrum, convolution, {windowing \& STFT}, mel}
\newcommand{\builtstations}{1}
\newcommand{\roadmap}[1]{%
  \begin{center}
  \begin{tikzpicture}[x=3.1cm, y=1cm,
      station/.style={rounded rectangle, draw=Muted, font=\scriptsize,
        inner sep=4pt, minimum height=5.5mm},
      done/.style={station, fill=GoodInk!25, draw=GoodInk},
      planned/.style={station, densely dashed, text=Muted},
      hop/.style={Muted, ->, shorten >=2pt, shorten <=2pt}]
    \foreach [count=\i] \name in \roadmapstations {
      \ifnum\i>\builtstations
        \node[planned] (s\i) at (\i, 0) {\name};
      \else
        \ifnum\i>#1
          \node[station] (s\i) at (\i, 0) {\name};
        \else
          \node[done] (s\i) at (\i, 0) {\name};
        \fi
      \fi
      \ifnum\i>1
        \pgfmathtruncatemacro{\prev}{\i - 1}
        \draw[hop] (s\prev.east) -- (s\i.west);
      \fi
    }
  \end{tikzpicture}
  \end{center}}


\title{From Air to Log-Mel\\[0.4em]
  \large A study guide, stitched from the manim-concepts road}
\author{manim-concepts}
\date{}

\begin{document}
\maketitle
\tableofcontents

\chapter*{How to read this guide}
\addcontentsline{toc}{chapter}{How to read this guide}
The objective is to read a log-mel spectrogram with understanding — the
picture a speech model is handed in place of sound — knowing what every
axis, every row and every shade of it measures, and what was thrown
away to make it. The road has four stations, one per video series of
the repo's signal-processing topic:

\roadmap{0}

\noindent This guide grows with the road. \emph{The spectrum} is
written: one column of the picture, built as a bank of detectors on
samples you can count. The dashed stations are planned and not yet
written — \emph{convolution} (the sliding weighted sum, and the theorem
that explains leakage), \emph{windowing, the short-time transform and
the spectrogram} (the columns side by side), and \emph{mel} (the rows
regrouped the way the ear groups them, with a log on top). Each is
written when it is asked for, and with it comes a transition page that
says why the next station needs exactly what the last one built.

The first chapter leans on three chapters of this repo's other guide,
\emph{The CTC Algorithm}: random variables (the weighted sum),
logarithms (multiplying is adding, and $\log_{10} 2$) and the counting
rules ($\binom{n}{k}$). Each use restates what it needs, so the chapter
reads on its own. Practice problems close every chapter — sized to
double as a diagnostic before you study digital audio — and worked
solutions live in the separate solutions manual, keyed by the same
problem numbers. Worked numbers are spliced from the repo's
verification anchors — for instance, the chapter's two-tone mix reads
$\anchor{019.R.readings}$ on its five detectors — never retyped.

% The retrieval manifest for this objective — the ordered \primitive calls
% ARE the remix map's spine (plan 012 R6); both build targets (guide and
% solutions manual) input this one file, so the stitching is single-sourced.
%
% Four chapters, one per main-line series of signal_processing/ (plan 019's
% road table). Only chapter 1 exists; the others are retrieved here when
% their series land (ADR 008), never as stubs:
%
%   1  spectrum                         built (plan 019)
%   2  convolution                      planned — series not built
%   3  windowing -> STFT -> spectrogram planned — series not built
%   4  mel                              planned — series not built
%
% No glue yet: a transition sits BETWEEN two chapters and is authored from
% the wiki edges its neighbours land, so glue-01.tex arrives with chapter
% 2, in the same change that adds its retrieval below and flips its
% station in macros.tex (\builtstations).

\primitive{spectrum}


\printbibliography[heading=bibintoc]

\end{document}
